% =============================================================================
% Step-5 supplement: unweighted vs integration on H2O training (8 archs).
% Compile alongside main.tex (shared references.bib, shared graphicspath).
% =============================================================================
\documentclass[11pt,letterpaper]{article}
\usepackage[margin=1in]{geometry}
\usepackage{graphicx, booktabs, amsmath, amssymb}
\usepackage[hidelinks]{hyperref}
\usepackage{natbib, xcolor, microtype, caption}
\captionsetup{font=small, labelfont=bf}
\graphicspath{
  {../step5_pretrain_origin_comparison/figures/}
  {../../notebooks/checkpoints_step5/unweighted/figures/}
  {../../notebooks/checkpoints_step5/integration/figures/}
}

\title{Step-5 Supplement: Unweighted vs.\ Integration Pretraining on H$_2$O}
\author{Alec P.\ Wills}
\date{\today}

\begin{document}
\maketitle

\section*{Run scope}
72 specs per origin = 8 architectures $\times$ 3 losses $\times$ 3 solvers,
trained on H$_2$O alone, evaluated on the W4-11 transfer set
$\{$CH$_4$, H$_2$, OH$\}$. Pretraining: 1000 steps under \texttt{x64} for each
of the 8 architectures, separately for each pretrain-loss-weighting choice.

\section{Headline numbers}

\begin{table}[h]
\centering
\caption{Step-5 headline numbers (verified against
\texttt{combined\_headline.json} and per-arch \texttt{pretrain\_metadata.json}).}
\small
\begin{tabular}{lrrr}
\toprule
metric & unweighted & integration & ratio \\
\midrule
best AE-MAE on H$_2$O (kcal/mol) & 0.001402 & \textbf{0.000118} & 11.9$\times$ tighter \\
best AE spec & deep\_dm/C/fixed\_j\_3 & deep\_dm\_attn/C/fixed\_j\_3 & adds attention \\
best $\rho$-RMSE on H$_2$O (e/bohr$^3$) & $4.41\!\times\!10^{-4}$ & $\mathbf{3.72\!\times\!10^{-4}}$ & 1.19$\times$ tighter \\
pretrain $F_x$ mean (8 archs) & $1.50\!\times\!10^{-4}$ & $\mathbf{3.52\!\times\!10^{-5}}$ & 4.26$\times$ tighter \\
pretrain $F_c$ mean (8 archs) & $4.80\!\times\!10^{-5}$ & $4.66\!\times\!10^{-5}$ & 1.03$\times$ (tied) \\
pretrained-only NN AE on H$_2$O (kcal/mol) & 27.2 & \textbf{17.1} & 1.59$\times$ tighter (37\%) \\
random-init NN AE on H$_2$O (kcal/mol) & 434.6 & 434.6 & tied \\
reduction vs PBE on H$_2$O ($\sim 7.03$ kcal/mol) & 5,000$\times$ & 59,500$\times$ & --- \\
\bottomrule
\end{tabular}
\end{table}

\noindent
The integration-pretrain origin produces an \textbf{11.9$\times$ tighter best
AE-MAE on H$_2$O} at the same training budget, before any downstream training
adjustments. The pretrained-only NN baseline is already 1.59$\times$ closer to
the answer (a 37\% AE-MAE reduction; the 1.59 ratio is \emph{not} a 59\%
reduction --- 1.59$\times$ tighter $\Leftrightarrow$ $(27.2-17.1)/27.2 = 37.1\%$
lower). That head start compounds through 250 task-loss steps to give the
12$\times$ best-AE improvement.

\section{Figures}

Figures generated by \texttt{notebooks/analysis/run\_step5\_origin\_comparison.py}
plus the in-notebook diagnostics. Per-step paths resolve via the shared
\texttt{\textbackslash{}graphicspath} preamble.

\begin{figure}[h!]
\centering
\includegraphics[width=\textwidth]{figs5_1_pretrain_loss.png}
\caption{Pretrain final $F_x$ (a) and $F_c$ (b) MSE per architecture, paired
bars (unweighted blue / integration green). $F_x$ ratios per arch: deep $4.74\times$,
deep\_attn $4.74\times$, deep\_dm $4.06\times$, deep\_dm\_attn $4.03\times$,
deep\_cusp $3.82\times$, deep\_cusp\_attn $4.51\times$, deep\_combined $2.95\times$,
deep\_combined\_attn $3.01\times$ --- range $\sim 3.0$ to $\sim 4.7$, mean $3.98\times$.
The earlier draft of this report claimed ``3.6$\times$ to 7.5$\times$'';
that was a transcription error and the correct range is given here.
$F_c$ is mixed: 5 of 8 architectures are \emph{worse} under integration on $F_c$
(deep\_dm, deep\_dm\_attn, deep\_cusp\_attn, deep\_combined, deep\_combined\_attn),
while 3 are tighter (deep, deep\_attn, deep\_cusp). The mean is essentially tied
because the 3 archs that benefit have the largest absolute $F_c$ baselines.}
\end{figure}

\begin{figure}[h!]
\centering
\includegraphics[width=\textwidth]{figs5_2_baseline_reduction.png}
\caption{H$_2$O AE-error baseline reduction. Trained-NN AE for all 72 specs per
origin; horizontal references are random-NN ($\sim$434 kcal/mol), pretrained-only
NN (27.2 vs.\ 17.1), PBE-on-H$_2$O ($\sim 7.03$ kcal/mol; this run's specific
def2-svp / grid\_level=1 H$_2$O configuration; \emph{not} a Karton 2011
literature value), and chemical accuracy 1 kcal/mol \citep{Pople1999}.}
\end{figure}

\begin{figure}[h!]
\centering
\includegraphics[width=\textwidth]{figs5_3_pareto_density_vs_AE.png}
\caption{Density-vs-energy plane. Each point is one spec; color = origin;
marker = architecture family. Integration green sits below-left of unweighted
blue --- a Pareto improvement on both axes simultaneously for most architectures.
Both clusters share the same density-RMSE floor at $\sim 3.7\times 10^{-4}$
e/bohr$^3$ (set by network expressive capacity / quadrature grid resolution,
not by pretrain origin).}
\end{figure}

\begin{figure}[h!]
\centering
\includegraphics[width=\textwidth]{figs5_4_arch_landscape.png}
\caption{Full AE-MAE landscape: 8 archs $\times$ 9 (loss $\times$ solver) cells
per origin, log color scale. Integration's lower-right (descriptor-rich attention
archs $\times$ C\_atomization\_plus\_grid loss) reaches deep green
(sub-millihartree), unmatched in the unweighted heatmap.}
\end{figure}

\begin{figure}[h!]
\centering
\includegraphics[width=\textwidth]{figs5_5_transfer.png}
\caption{Transfer-mol AE per architecture and origin. PBE references are
hardcoded for our specific def2-svp / grid\_level=1 configuration:
CH$_4$ 0.91, H$_2$ 7.31, OH 1.84 kcal/mol. These are \emph{not} W4-11
literature values --- \citet{KartonDaonMartin2011} tabulates CCSD(T)-quality
reference atomization energies, not PBE errors. No spec on any architecture
or origin reaches chemical accuracy on any of CH$_4$, H$_2$, OH:
H$_2$O-only training is insufficient for the closed-shell light-hydride
class. OH is additionally open-shell ($^2\Pi$, $S{=}1/2$); the pretrain
\emph{atom} set $\{$H, He, O, N$\}$ does include open-shell atoms, so the
mismatch is on the molecular spin-density profile, not on atomic
open-shell character.}
\end{figure}

\section{In-notebook diagnostic figures}

The step-5 notebook (\texttt{notebooks/\_build\_step5\_notebook.py})
produces the following diagnostics directly. They cover pretrain
dynamics (\texttt{pretrain\_losses.png}, \texttt{pretrain\_parity.png}),
training-phase loss curves, auxiliary-loss component breakdowns
(\texttt{aux\_components\_deep\_combined.png}), an exploratory
loss-balancing pilot (\texttt{balancing\_aux\_comparison.png},
\texttt{bal\_loss\_plot\_vxc.png} --- the latter motivated step 6's
$V_{xc}$-aware losses but is not part of step 5's production
$\{A,B,C\}$ loss family), GradNorm task-weight evolution, SCF
convergence aggregates, density-matrix and grid-density diagnostics,
$F_x(s)$ drift on the CH$_4$ molecular grid, transfer aggregate and
per-mol breakdowns, and a feature-impact analysis. Captions describe
each plot's semantics; the diagnostics are for completeness, not the
headline results (\S1).

\begin{figure}[h!]
\centering
\includegraphics[width=0.92\textwidth]{pretrain_losses.png}
\caption{Pretrain loss curves per architecture (in-notebook).}
\end{figure}

\begin{figure}[h!]
\centering
\includegraphics[width=0.92\textwidth]{pretrain_parity.png}
\caption{Pretrain parity plots per architecture (in-notebook).}
\end{figure}

\begin{figure}[h!]
\centering
\includegraphics[width=0.92\textwidth]{training_losses.png}
\caption{Training-phase loss curves per architecture/loss/solver (in-notebook).}
\end{figure}

\begin{figure}[h!]
\centering
\includegraphics[width=0.6\textwidth]{aux_components_deep_combined.png}
\caption{Auxiliary loss components: deep\_combined arch.}
\end{figure}

\begin{figure}[h!]
\centering
\includegraphics[width=0.92\textwidth]{balancing_aux_comparison.png}
\caption{Loss-balancing comparison aux plot (in-notebook).}
\end{figure}

\begin{figure}[h!]
\centering
\includegraphics[width=0.6\textwidth]{bal_loss_plot_vxc.png}
\caption{Loss-balancing experiment loss curves on the
\texttt{deep\_combined} architecture (in-notebook). This is the
exploratory pilot run that motivated the step-6 $V_{xc}$-aware loss
families; step 5's actual production losses
$\{A, B, C\}$ are exchange-and-density only and do not include a
$V_{xc}$ residual.}
\end{figure}

\begin{figure}[h!]
\centering
\includegraphics[width=0.92\textwidth]{gradnorm_weight_evolution.png}
\caption{GradNorm task-weight evolution \citep{ChenGradNorm2018}.}
\end{figure}

\begin{figure}[h!]
\centering
\includegraphics[width=0.92\textwidth]{scf_comparison_ae.png}
\caption{SCF-mode AE comparison.}
\end{figure}

\begin{figure}[h!]
\centering
\includegraphics[width=0.6\textwidth]{scf_convergence.png}
\caption{SCF convergence aggregate.}
\end{figure}

\begin{figure}[h!]
\centering
\includegraphics[width=0.92\textwidth]{dm_heatmaps_scf.png}
\caption{DM heatmaps after final SCF.}
\end{figure}

\begin{figure}[h!]
\centering
\includegraphics[width=0.92\textwidth]{grid_density_scf.png}
\caption{Grid-density profile after final SCF.}
\end{figure}

\begin{figure}[h!]
\centering
\includegraphics[width=0.6\textwidth]{fx_drift_ch4.png}
\caption{$F_x(s)$ drift sampled on the CH$_4$ molecular grid.}
\end{figure}

\begin{figure}[h!]
\centering
\includegraphics[width=0.92\textwidth]{transfer_aggregate.png}
\caption{Transfer aggregate plot (in-notebook).}
\end{figure}

\begin{figure}[h!]
\centering
\includegraphics[width=0.6\textwidth]{transfer_CH4.png}
\caption{CH$_4$ transfer detail (in-notebook).}
\end{figure}

\begin{figure}[h!]
\centering
\includegraphics[width=0.6\textwidth]{transfer_H2.png}
\caption{H$_2$ transfer detail (in-notebook).}
\end{figure}

\begin{figure}[h!]
\centering
\includegraphics[width=0.6\textwidth]{transfer_OH.png}
\caption{OH transfer detail (in-notebook).}
\end{figure}

\begin{figure}[h!]
\centering
\includegraphics[width=0.92\textwidth]{feature_impact_scf.png}
\caption{Feature-impact analysis (in-notebook).}
\end{figure}

\section{Conclusions and recommendations}

\begin{enumerate}
\item Use the integration-pretrain origin for any subsequent training that targets AE
  accuracy. The $F_x$ pretrain is tighter on every architecture (8/8 in step 5);
  the pretrained-only NN baseline is 1.6$\times$ tighter; the downstream best AE is
  $\sim$12$\times$ tighter.
\item For density-quality priority, integration + C\_atomization\_plus\_grid +
  descriptor-rich arch (deep\_dm\_attn / deep\_cusp\_attn / deep\_combined\_attn) +
  oneshot or fixed\_j\_3 solver gives the best Pareto point on the density-vs-AE
  plane ($\sim 3.7 \times 10^{-4}$ e/bohr$^3$ at sub-mHa AE).
\item Do not expect transfer-set generalization gains from any pretrain-weighting
  variant. Both origins give 5--21 kcal/mol AE on $\{$CH$_4$, H$_2$, OH$\}$ ---
  comparable to or worse than PBE on each. Scale the training-molecule set
  (\citet{BehlerParrinello2007}, by analogy for NN XC).
\item $F_c$ pretrain quality is comparable across both origins, and the implementation
  already does the right thing: \texttt{xcquinox/alec/pretrain.py:69--77} builds
  $w_x$ from $|\rho \varepsilon_x^{\text{LDA}}|$ and $w_c$ from
  $|\rho \varepsilon_c^{\text{PW92}}|$ \citep{PerdewWang1992} separately. An earlier
  recommendation to ``switch the F$_c$ weight to PW92'' was based on a misreading of
  the code; the fix is already in place.
\end{enumerate}

\bibliographystyle{plainnat}
\bibliography{references}
\end{document}
